\documentclass[../Main.tex]{subfiles}

\begin{document}
\section{Formulations}
\begin{definition}{Linear program}
    A \underline{linear program} is an optimisation problem where the objective function and all constraints are linear in $\vec{x}$.
\end{definition}
The \underline{general form} is:
\begin{align*}
    \text{Maximise } & \vec{c} \cdot \vec{x} \\
    \text{Subject to } & \vec{a_i} \cdot \vec{x} \geq \vec{b_{i}},~~i \in M_1 \\
    & \vec{a_i} \cdot \vec{x} \leq \vec{b_{i}},~~i \in M_2 \\
    & \vec{a_i} \cdot \vec{x} = \vec{b_{i}},~~i \in M_3 \\
    & x_j \geq 0,~~j \in N_1 \\
    & x_j \leq 0,~~j \in N_2.
\end{align*}
The \underline{standard form} is:
\begin{align*}
    \text{Maximise } & \vec{c} \cdot \vec{x} \\
    \text{Subject to } & A\vec{x} = \vec{b} \\
    & \vec{x} \geq \vec{0}.
\end{align*}
Any problem in general form can be written in standard form using slack variables and by expressing components as $x_j = x_j^+ - x_j^-$.
\section{Extreme Points}
\begin{definition}{Extreme point}
    A point $\vec{x}$ of a convex set $\Xset \subseteq \R^n$ is an \underline{extreme point} if $\vec{x}$ cannot be written as:
    \begin{equation*}
        \vec{x} = \lambda \vec{y} + (1 - \lambda) \vec{z}
    \end{equation*}
    for $\vec{y} \neq \vec{z}$ and $\lambda \in (0, 1)$.
\end{definition}
\begin{theorem}
    The maximum of a convex function over a convex set is realised at an extreme point of this set.
    \label{thmMaxAtExtreme}
\end{theorem}
\begin{proof}
    A simple proof by contradiction.
\end{proof}

For a matrix $A \in \R^{m \times n}$, let $\vec{a_{1}}, \cdots \vec{a_{m}}$ be the rows of the matrix and $\vec{A_{1}}, \cdots, \vec{A_{n}}$ be the columns.

Now consider a linear programming problem in standard form and assume:
\begin{enumerate}
    \item $A$ has linearly independent rows, so the rank is $m$ (if not we would have redundant constraints we could remove)
    \item Every set of $m$ columns is linearly independent ($n > m$)
\end{enumerate}
\begin{definition}{Basic solution}
    A point $\vec{x}$ that satisfies $A\vec{x} = \vec{b}$ and has at most $m$ non-zero elements is a \underline{basic solution} of the linear programming problem.
\end{definition}
\begin{definition}{Basic feasible solution}
    A basic solution $\vec{x}$ that satisfies $\vec{x} \geq \vec0$ is a \underline{basic feasible solution}.
\end{definition}
We now have a third assumption:
\begin{enumerate}
    \setcounter{enumi}{2}
    \item Every basic feasible solution is non-degenerate (has exactly $m$ non-zero entries).
\end{enumerate}
We can now detail a simple algorithm to find basic solutions:
\begin{enumerate}
    \item Select indices $B(1), \cdots, B(m) \in \{1, \cdots, n\}$.
    \item For each index $B(i)$, set the $i$th column of a matrix $B$ to be $\vec{A_{B(i)}}$
    \item By assumption 2, $B$ is invertible: set the $B(i)$th element of $\vec{x}$ to $[B^{-1}\vec{b}]_i$
    \item Set the other entries of $\vec{x}$ to zero.
\end{enumerate}
Then to find basic feasible solutions, we simply check each basic solution.
\begin{theorem}[Basic feasible solutions are extreme points]
    A point $\vec{x}$ is a basic feasible solution if and only if it is an extreme point of the feasible set:
    \begin{equation*}
        \Xset(\vec{b}) = \subsetselect{\vec{x} \in \R^n}{A\vec{x} = \vec{b},~\vec{x} \geq \vec{0}}
    \end{equation*}
    \label{thmBFSIsExtreme}
\end{theorem}
\begin{proof}
    \begin{proofdirection}{$\Rightarrow$}{Assume that $\vec{x}$ is a BFS}
        Proof by contradiction: let $\vec{x} = \lambda \vec{y} + (1-\lambda)\vec{z}$. Then show that these must have the same zero entries (because they are all feasible, $\vec{x} \geq \vec0$) and then are determined by uniqueness of solution to $B\vec{v} = \vec{b}$.
    \end{proofdirection}
    \begin{proofdirection}{$\Leftarrow$}{Assume that $\vec{x}$ is an extreme point}
        Proof by contradiction: let $\vec{x}$ be extreme but not a BFS, so has $r > m$ non-zero entries. Label the indices $i_1, \cdots, i_r$ and take the columns of $A$ at these indices. Take a linear combination of these columns and use the coefficients to form a new vector such that $A \vec{w} = \vec{0}$. Then for small enough $k$ such that $\vec{x} \pm k\vec{w} \geq 0$, we have a convex combination to make $\vec{x}$ which contradicts the extremeness.
    \end{proofdirection}
\end{proof}
\section{Duality in Linear Programs}
\begin{theorem}
    If a linear program is feasible and bounded then it satisfies strong duality.
    \label{thmLinearProgStrongDual}
\end{theorem}
\begin{proof}
    The proof is by finding a supporting hyperplane.
\end{proof}
For a linear program in standard form:
\begin{align*}
    \text{Maximise } & \vec{c} \cdot \vec{x} \\
    \text{Subject to } & A\vec{x} = \vec{b} \\
    & \vec{x} \geq \vec{0}
\end{align*}
its dual is:
\begin{align*}
    \text{Maximise} & \vec{b} \cdot \vec{\lambda} \\
    \text{Subject to } & A^T\vec{\lambda} \leq \vec{c}.
\end{align*}

For a linear program in general form:
\begin{align*}
    \text{Maximise } & \vec{c} \cdot \vec{x} \\
    \text{Subject to } & \vec{a_i} \cdot \vec{x} \geq \vec{b_{i}},~~i \in M_1 \\
    & \vec{a_i} \cdot \vec{x} \leq \vec{b_{i}},~~i \in M_2 \\
    & \vec{a_i} \cdot \vec{x} = \vec{b_{i}},~~i \in M_3 \\
    & x_j \geq 0,~~j \in N_1 \\
    & x_j \leq 0,~~j \in N_2 \\
    & x_j \text{ free},j \in N_3
\end{align*}
its dual is:
\begin{align*}
    \text{Maximise } & \vec{b} \cdot \vec{\lambda} \\
    \text{Subject to } & \lambda_i \geq 0,~~i \in M_1 \\
    & \lambda_i \leq 0,~~i \in M_2 \\
    & \lambda_i \text{ free},~~i \in M_3 \\
    & \vec{\lambda}^T \vec{A_{j}} \leq c_j,~~j \in N_1 \\
    & \vec{\lambda}^T \vec{A_{j}} \geq c_j,~~j \in N_2 \\
    & \vec{\lambda}^T \vec{A_{j}} = c_j,~~j \in N_3 \\
\end{align*}
\begin{theorem}
    Let $\vec{x}$ and $\vec{\lambda}$ be feasible points in the primal and dual problems, respectively. The pair $(\vec{x}, \vec{\lambda})$ is optimal if and only if:
    \begin{align*}
        \lambda_i (\vec{a_{i}}^T \vec{x} - b_i) &= 0 ~\forall i \\
        (c_i - \vec{\lambda}^T \vec{A_{j}})x_j &= 0 ~\forall j
    \end{align*}
    \label{thmLinearProgOptimality}
\end{theorem}
\begin{remark}
    Linear programs are special in that complementary slackness is sufficient for optimality.
\end{remark}
\begin{proof}
    Define:
    \begin{align*}
        u_i &= \lambda_i (\vec{a_{i}}^T \vec{x} - b_i) \\
        v_j &= (c_i - \vec{\lambda}^T \vec{A_{j}})x_j
    \end{align*}
    These are non-negative. Summing this all together gives $\vec{c} \cdot \vec{x} - \vec{\lambda} \cdot \vec{b}$.

    If complementary slackness holds then all the $u_i$ and $v_j$ are zero so we have strong duality. This is reversible logic because they are all non-negative.
\end{proof}
\section{Simplex Algorithm}
\begin{definition}{Reduced cost vector}
    Given a set of indices $B(1), \cdots, B(m)$, the \underline{reduced cost vector} $\bvec{c}$ is defined:
    \begin{equation*}
        \bar{c}_j = c_j - \vec{c_B}^T B^{-1} \vec{A_{j}}
    \end{equation*}
    where $\vec{c_{B}}$ is the vector with only elements $B(1), \cdots, B(m)$ and $B$ is the matrix of the corresponding columns of $A$ (as before).
\end{definition}
A basic feasible solution is optimal if and only if $\bvec{c} \geq \vec0$.

The simplex algorithm starts at a basic feasible solution and moves in the direction of decreasing cost. It considers a simplex Tableau, and the following algorithm.
\begin{enumerate}
    \item Setup: add slack variables as required to get to standard form.
    \item Start the simplex method with a basic feasible solution (such as by setting all non-slack variables to 0 if this works).
    \item Fill in the simplex tableau with:
    \begin{enumerate}
        \item the non-zero entries of $\vec{x}$ in the left column,
        \item the corresponding rows of $B^{-1} A$ in the body of the table,
        \item the reduced costs in the bottom row: $\bar{c}_i$,
        \item the value $-\vec{c_B} \cdot \vec{x_B}$, to minimise.
    \end{enumerate}
    \item Identify the pivot column: if all entries in the bottom row ($\bvec{c}$) are non-negative, stop. If not, select the smallest negative entry.
    \item Identify the pivot element: for the pivot column, divide the elements in the first row by the corresponding elements in the pivot column. Choose the row corresponding to the smallest ratio.
    \item Perform row operations: the $j$th column should be the identity column so scale the $i$th row and add it to the other rows as required.
    \item Return to step 3 and continue the algorithm.
\end{enumerate}
\end{document}